% Part: lambda-calculus
% Chapter: church-rosser
% Section: beta-eta-reduction

\documentclass[../../../include/open-logic-section]{subfiles}

\begin{document}

\olfileid{lam}{cr}{be}

\olsection{$\beta\eta$-归约}

Church（邱奇）--Rosser（罗瑟）性质对 $\beta\eta$-归约（$\bered$）成立。

\begin{lem}\ollabel{lem:one-par}
  若 $M \beredone M'$，则 $M \beredpar M'$。
\end{lem}

\begin{proof}
  对 $M \beredone M'$ 的推导过程进行归纳。如果 $M \bredone M'$ 是通过 $\eta$-转换（即 \olref[syn][eta]{defn:beredone}）得到的，我们使用 \olref[pbe]{thm:refl}。其余情况同 \olref[b]{lem:one-par} 的证明。
\end{proof}

\begin{lem}\ollabel{lem:par-red}
  若 $M \beredpar M'$，则 $M \bered M'$。
\end{lem}

\begin{proof}
  对 $M \beredpar M'$ 的推导过程进行归纳。

  如果最后一条规则是 \olref[pbe]{defn:beredpar5}，那么 $M$ 形如 $\lambd[x][Nx]$ 且 $M'$ 为 $N'$，其中 $x \notin FV(N)$ 且 $N \beredpar N'$。因此，我们可以先通过 $\eta$-转换将 $\lambd[x][Nx]$ 归约为 $N$，随后进行由归纳假设保证的一系列 $\beredone$ 步骤，从而得出 $N \bered N'$。
\end{proof}

\begin{lem}\ollabel{lem:str}
  $\bered$ 是包含 $\beredpar$ 的最小传递关系。
\end{lem}

\begin{proof}
  同 \olref[b]{lem:str} 的证明。
\end{proof}

\begin{thm}\ollabel{thm:cr}
  $\bered$ 满足 Church--Rosser 性质。
\end{thm}

\begin{proof}
  由 \olref[dap]{thm:str}、\olref[pbe]{thm:cr} 以及 \olref{lem:str} 可得。
\end{proof}

\end{document}